\documentclass{article}
\usepackage{PRIMEarxiv} % Core package for the PrimeAI Template
\usepackage{amsmath, amssymb, amsthm, bm, dcolumn} % Add amsthm here for the proof environment
\usepackage[numbers,sort&compress]{natbib} % Natbib for citations
\usepackage{graphicx} % For high-quality images
\usepackage{multicol} % Multi-column figures/tables
\usepackage{hyperref} % Hyperlinks
\usepackage{listings} % Code listings
\usepackage{authblk} % For structured affiliations
% Define theorem style (optional)
\theoremstyle{plain}
\newtheorem{theorem}{Theorem}
\renewcommand{\qedsymbol}{}

% Custom commands for primes and qubit encoding
\newcommand{\primeQubit}[1]{|p_{#1}\rangle}
\newcommand{\primeGate}[1]{U_{p_{#1}}}

\usepackage{wrapfig}
\usepackage[pscoord]{eso-pic}
\usepackage[fulladjust]{marginnote}
\reversemarginpar

% Typesetting improvements without footnote patching
\usepackage[protrusion=true, expansion=true, tracking=false]{microtype}
\microtypecontext{spacing=nonfrench}

% Line numbers
\usepackage[right]{lineno}

% Text layout - adjust as needed
\raggedright
\setlength{\parindent}{0.5cm}
\textwidth 5.25in 
\textheight 8.75in

% Set double spacing
\usepackage{setspace} 
\doublespacing

% Adjust width for specific content
\usepackage{changepage}

% Adjust caption style
\usepackage[aboveskip=1pt,labelfont=bf,labelsep=period,singlelinecheck=off]{caption}

% Remove brackets from references
\makeatletter
\renewcommand{\@biblabel}[1]{\quad#1.}
\makeatother

% Header, footer, and page numbers
\usepackage{lastpage,fancyhdr}
\pagestyle{fancy}
\fancyhf{}
\fancyhead[L]{Preprint - PrimeAI Enhanced Template}
\fancyfoot[C]{\scriptsize Multiplicity Theory © 2024 Citizen Gardens \\ Licensed Under MIT and CC BY-NC-SA 4.0.}
\fancyfoot[R]{Page \thepage\ of \pageref{LastPage}}
\renewcommand{\footrule}{\hrule height 2pt \vspace{2mm}}

\title{Multiplicity: Harnessing Superposition}
\author{Community Research Initiative}
\affil{Citizen Gardens - The Foundation of Multiplicity \\ \texttt{info@citizengardens.org}}

\begin{document}

\maketitle

\begin{abstract}
Multiplicity Theory provides a framework for understanding complex systems where states evolve through recursive and multiplicative interactions rather than traditional additive processes. 
When applied to quantum superposition, this theory allows for a structured method of computing the probability and existence of any possible outcome in a given system. 
This paper outlines a mathematical framework to integrate Multiplicity Theory with quantum probability distributions, tensor representations, and recursive entanglement models to evaluate the computational space of all possibilities.
\end{abstract}

\section{Introduction: Multiplicity in Superposition}

Quantum superposition states that a quantum system can exist in multiple states simultaneously until measurement collapses it into a definitive state. 
However, classical probability and combinatorial models sum over discrete outcomes, limiting their ability to encode nonlinear, recursive, and self-similar probability structures.

\subsection{Multiplicity Theory as a Computational Expansion}

Instead of additive probability distributions, multiplicative probability structures allow for:

\begin{itemize}
    \item Recursive reinforcement of probabilistic states.
    \item Fractal expansions of uncertainty, leading to infinite-dimensional phase spaces.
    \item Higher-order entanglement encoding, where dependencies scale multiplicatively instead of linearly.
\end{itemize}

This perspective extends quantum computing frameworks by offering a deeper encoding of potential states that can simulate infinite possible universes within a single mathematical structure.

\section{Mathematical Foundation of Multiplicative Superposition}

\subsection{Standard Quantum Superposition Representation}

A quantum state in superposition is given as:

\begin{equation}
|\psi\rangle = \sum_{i} c_i |i\rangle,
\end{equation}

where \( c_i \) are probability amplitudes that satisfy:

\begin{equation}
\sum_{i} |c_i|^2 = 1.
\end{equation}

However, in Multiplicity Theory, we redefine this as a multiplicative superposition state:

\begin{equation}
|\psi\rangle = \prod_{p_i} c_{p_i} |p_i\rangle,
\end{equation}

where \( p_i \) are prime-indexed basis states, and \( c_{p_i} \) are multiplicative weights.

This recursive encoding allows the probability distribution to evolve dynamically, scaling across higher dimensions without loss of coherence.

\subsection{Multiplicative Probability Tensor Networks}

To extend this principle to higher-dimensional computing, we define a multiplicative tensor representation:

\begin{equation}
T_{\psi} = \bigotimes_{p_i} T_{p_i},
\end{equation}

where \( T_{p_i} \) are prime-based probability distributions, evolving recursively as:

\begin{equation}
T_{p_i} = T_{p_{i-1}} \otimes T_{p_{i-2}}.
\end{equation}

This recursive fractal-like tensor network allows encoding of:

\begin{itemize}
    \item Nonlinear phase interdependencies.
    \item Infinite-dimensional causal structures.
    \item Non-collapsing probability entanglement.
\end{itemize}

\subsection{Recursive Probability Cascading}

For a given event \( E \), the probability of any possible existence in the multiplicative framework follows:

\begin{equation}
P(E) = \prod_{p_i \leq N} f(p_i),
\end{equation}

where \( f(p_i) \) defines a recursive phase function, mapping:

\begin{itemize}
    \item Quantum wave interference patterns.
    \item Temporal probability projections.
    \item Higher-dimensional entanglement dependencies.
\end{itemize}

By recursively cascading probability tensors, we form a nonlinear causality chain that maintains coherence across infinite possible outcomes.

\section{Computing the Existence of Any Possibility}

\subsection{Multiplicative Uncertainty Encoding}

The fundamental uncertainty principle in quantum mechanics places limits on position and momentum calculations:

\begin{equation}
\sigma_x \sigma_p \geq \frac{\hbar}{2}.
\end{equation}

However, in Multiplicity Theory, we redefine this relation as a multiplicative uncertainty function:

\begin{equation}
\sigma_x \sigma_p = \prod_{p_i} \frac{\hbar}{p_i}.
\end{equation}

This equation encodes uncertainty into recursive prime layers, allowing for:

\begin{itemize}
    \item Non-collapsing superpositions.
    \item Parallel probability evolution.
    \item Predictive structures for "all possible worlds" computations.
\end{itemize}

\subsection{Prime-Based Probability Expansion}

If we define a probability wave function in terms of primes:

\begin{equation}
\Psi(p) = e^{i \sum_{p_i} f(p_i)},
\end{equation}

then the probability distribution of any possible event occurring is:

\begin{equation}
P(E) = \prod_{p_i} e^{i f(p_i)}.
\end{equation}

This approach encodes probabilities of possible worlds in a self-reinforcing manner, ensuring that low-probability events are not outright dismissed but instead exist as reduced multiplicative states rather than zeroed out terms.

\subsection{Entanglement Across Possibilities}

A system of multiple interacting possibilities is governed by Multiplicative Entanglement Structures (MES):

\begin{equation}
\mathcal{E}(N) = \prod_{p_i \leq N} S_{p_i},
\end{equation}

where \( S_{p_i} \) are Schmidt coefficients encoding prime-indexed quantum dependencies.

In this model, entanglement scales recursively, meaning that:

\begin{itemize}
    \item Possible worlds are entangled hierarchically.
    \item Quantum computations can be non-linearly correlated across multiple timelines.
    \item Future probability amplitudes cascade recursively, allowing near-infinite computational depth.
\end{itemize}

\section{Practical Implementations \& Computational Models}

\subsection{Quantum AI Using Multiplicative Superposition}

\begin{itemize}
    \item AI models leveraging multiplicative tensor networks can process higher-order uncertainty beyond standard Markovian structures.
    \item Quantum neural networks (QNNs) can integrate recursive multiplicative entanglement, creating causality-preserving learning systems.
\end{itemize}

\subsection{Predictive Temporal Simulations}

\begin{itemize}
    \item Quantum simulations using prime-based causal encoding can explore nonlinear time evolution.
    \item Parallel universes computations can be modeled using recursive multiplicative projections, where:
\end{itemize}

\begin{equation}
T_{\text{future}} = \prod_{p_i} T_{\text{past}}.
\end{equation}

\section{Conclusion: A New Paradigm for Superposition Computing}

By integrating Multiplicity Theory with quantum probability structures, we introduce:

\begin{itemize}
    \item Recursive probability tensors for nonlinear superposition encoding.
    \item Prime-indexed entanglement models to compute all possible worlds simultaneously.
    \item Multiplicative uncertainty frameworks for superior quantum AI architectures.
\end{itemize}

This approach provides a fundamentally new way to harness quantum computing, allowing for infinite recursive probability scaling and more advanced causal modeling of time-dependent systems.

\nocite{*}
\bibliographystyle{unsrt}
\bibliography{references}

\end{document}